% Generated draft: requires mathematical and source review.
\documentclass[11pt]{article}
\usepackage[margin=25mm]{geometry}
\usepackage{fontspec}
\setmainfont{Latin Modern Roman}
\usepackage{amsmath,amssymb,mathtools,unicode-math}
\setmathfont{Latin Modern Math}
\usepackage{enumitem,needspace,xurl,hyperref}
\setlength{\parindent}{0pt}
\setlength{\parskip}{5pt}
\setlength{\emergencystretch}{3em}
\newcommand{\N}{\mathbb{N}}
\newcommand{\Z}{\mathbb{Z}}
\newcommand{\Q}{\mathbb{Q}}
\newcommand{\R}{\mathbb{R}}
\newcommand{\C}{\mathbb{C}}
\newcommand{\F}{\mathbb{F}}
\newcommand{\A}{\mathbb{A}}
\newcommand{\PP}{\mathbb{P}}
\newcommand{\E}{\mathbb{E}}
\newcommand{\eps}{\varepsilon}
\newcommand{\dd}{\,\mathrm{d}}
\begin{document}
\section*{UnsolvedMath ID: 30003436}
\textbf{Source problem code:} catalogue entry
\subsection{Definitions and mathematical statement}
For integers \(n\geq 1\) and \(0\leq k<n\), let
\[
G_{k,n}=\operatorname{Gr}_{\R}(k+1,n+1)
\]
be the Grassmannian parametrizing projective \(k\)-dimensional linear subspaces of \(\mathbb{RP}^{n}\). The source denotes its expected degree by \(\delta_{k,n}\), but the supplied record does not define the random incidence experiment or its probability distribution. Consequently, the normalization and precise probabilistic definition of \(\delta_{k,n}\) require confirmation from the cited report.

In the real projective cubic-surface interpretation suggested by the context, let \(F\in\R[x_0,x_1,x_2,x_3]\) be a nonzero homogeneous polynomial of degree three, and set
\[
X_F=\{[x]\in\mathbb{RP}^{3}:F(x)=0\},\qquad
\mathcal{L}(F)=\{L\in G_{1,3}:L\subseteq X_F\}.
\]
A random cubic surface requires a specified probability law for \(F\), or for its class under multiplication by nonzero real scalars. The source question does not specify this law or whether smoothness is assumed. The reference to the Kostlan ensemble in the supplied literature assessment does not establish that this ensemble is required by the question.

The original source record asks the following three questions:
\begin{enumerate}
\item Is there an explicit formula in special functions for \(\delta_{1,3}\)?
\item For each fixed admissible integer \(k\), what is the asymptotic behavior of \(\delta_{k,n}\) as \(n\to\infty\)? The record does not specify the desired precision of the asymptotic description.
\item Is there a convex body naturally associated with the lines on a random cubic surface, analogous to the convex-geometric objects used to study expected degrees of real Grassmannians?
\end{enumerate}
The current statement in the registry selects the third question. Here a convex body means a compact convex subset of a finite-dimensional real vector space with nonempty interior relative to its affine hull. This convention alone does not make the question a precise existence assertion: the record does not specify the ambient vector space, the required association with \(\mathcal{L}(F)\) or its distribution, the meaning of naturality, or the property that realizes the proposed analogy. These requirements remain unresolved; no particular construction or identity is asserted.

The supplied dated curation labels the overall problem partially solved. It reports the Kostlan expected number of real lines as \(6\sqrt{2}-3\), an explicit special-function or elliptic-integral formula for \(\delta_{1,3}\), and proved fixed-\(k\) asymptotics, while describing the convex-body question as open. These are supplied literature assessments, not independently verified mathematical or current-status claims.
\subsection{Short English statement}
The current question asks for a natural convex body associated with lines on a random cubic surface, in analogy with convex geometry used for expected degrees of real Grassmannians. A precise existence statement still requires the random-surface model, the intended association, and a defining property of the convex body. The original record also asks about an explicit formula and asymptotics for expected degrees; supplied dated curation reports progress on those parts.
\subsection{Sources}
\begin{itemize}
\item Local UnsolvedMath registry, permanent ID \(30003436\), source problem code \texttt{catalogue entry}, title \emph{Expected Degrees of Random Real Grassmannians}, Oberwolfach Reports Open Problems collection. No individual proposer is supplied; the proposed year is recorded as 2017. The registry labels the statement \texttt{corrected\_verified}, explaining that notation, typography, and sentence structure were repaired while preserving mathematical content. This is a supplied editorial qualification, not independent mathematical verification.
\item Supplied report citation: \emph{Real Algebraic Geometry With a View Toward Moment Problems and Optimization} (2018), Oberwolfach Reports, pp.\ 771--862, DOI \texttt{10.4171/owr/2017/14}. \url{https://doi.org/10.4171/owr/2017/14}. The supplied record attributes the problem to workshop open-problem discussions but gives no precise page or named contributor.
\item Basu--Lerario--Lundberg--Peterson, supplied reference label: random cubic lines. \url{https://arxiv.org/abs/1610.01205}.
\item Lerario--Mathis, supplied reference label: probabilistic Schubert asymptotics. \url{https://arxiv.org/abs/1912.08291}.
\item Ait El Manssour--Belotti--Meroni, supplied reference label: random cubic surfaces. \url{https://doi.org/10.1007/s40598-021-00182-y}.
\item Literature assessment checked at \texttt{2026-08-21}: supplied dated curation derived from the OWR record, not independently verified. It labels the problem partially solved, reports the Kostlan expectation \(6\sqrt{2}-3\), an explicit formula for \(\delta_{1,3}\), and fixed-\(k\) asymptotics, and describes the convex-body question as open. The record explicitly requests independent verification of time-sensitive status and cited claims.
\end{itemize}
\end{document}
